% FAC-EN-U0058
% Source: fac_body.tex lines 2079--2130.
% Authority: PDF physical pp. 57--58; printed pp. 252--253 (no. 60 complete).
% English source-first translation; printed upper-row L symbols corrected and logged.

\phantomsection\label{fac:III:2:no:60}
\textbf{60. The case of coherent algebraic sheaves.}

Let us first prove a preliminary result.

\phantomsection\label{fac:III:2:no:60:prop:9}
\textsc{Proposition 9.} \textit{Let $\mathcal{L}$ be an algebraic sheaf
on $X$ that is a direct sum of finitely many sheaves $\mathcal{O}(n_i)$.
Then $\Gamma(\mathcal{L})$
\hyperref[fac:III:2:no:56:def:TF]{satisfies condition (TF)}, and
$\hyperref[fac:III:2:no:59:def:beta]{\beta}:
\mathcal{A}(\Gamma(\mathcal{L}))\to\mathcal{L}$ is bijective.}

We reduce at once to $\mathcal{L}=\mathcal{O}(n)$, and then to
$\mathcal{L}=\mathcal{O}$. In this case we know that
$\Gamma(\mathcal{O}(p))=S_p$ for $p\geq0$, and hence
$S\subset\Gamma(\mathcal{O})$, with the quotient belonging to
$\mathcal{C}$. It follows first that $\Gamma(\mathcal{O})$
\hyperref[fac:III:2:no:56:def:TF]{satisfies (TF)}, and then that
\[
\mathcal{A}(\Gamma(\mathcal{O}))=
\mathcal{A}(S)=\mathcal{O},
\]
which proves the proposition.

(Observe that $\Gamma(\mathcal{O})=S$ if $r\geq1$; if $r=0$, however,
$\Gamma(\mathcal{O})$ is not even a finitely generated $S$-module.)

\phantomsection\label{fac:III:2:no:60:theorem:2}
\textsc{Theorem 2.} \textit{For every coherent algebraic sheaf
$\mathcal{F}$ on $X$, there is a graded $S$-module $M$
\hyperref[fac:III:2:no:56:def:TF]{satisfying condition (TF)} such that
$\mathcal{A}(M)$ is isomorphic to
$\mathcal{F}$.}

By the \hyperref[fac:III:2:no:55:corollary]{corollary to Theorem~1},
there is an exact sequence of algebraic sheaves
\phantomsection\label{fac:III:2:no:60:formula:sheaf-presentation}
\[
\mathcal{L}^1\xrightarrow{\varphi}\mathcal{L}^0
\longrightarrow\mathcal{F}\longrightarrow0,
\]
where $\mathcal{L}^1$ and $\mathcal{L}^0$ satisfy the hypotheses of the
\hyperref[fac:III:2:no:60:prop:9]{preceding proposition}. Let $M$ be the
cokernel of the homomorphism
\phantomsection\label{fac:III:2:no:60:formula:Gamma-phi}
\[
\Gamma(\varphi):\Gamma(\mathcal{L}^1)
\longrightarrow\Gamma(\mathcal{L}^0).
\]
By \hyperref[fac:III:2:no:60:prop:9]{Proposition~9}, $M$
\hyperref[fac:III:2:no:56:def:TF]{satisfies condition (TF)}. Applying the
functor $\mathcal{A}$ to the exact sequence
\phantomsection\label{fac:III:2:no:60:formula:module-cokernel}
\[
\Gamma(\mathcal{L}^1)\longrightarrow
\Gamma(\mathcal{L}^0)\longrightarrow M\longrightarrow0,
\]
we obtain the exact sequence
\phantomsection\label{fac:III:2:no:60:formula:sheafified-cokernel}
\[
\mathcal{A}(\Gamma(\mathcal{L}^1))\longrightarrow
\mathcal{A}(\Gamma(\mathcal{L}^0))\longrightarrow
\mathcal{A}(M)\longrightarrow0.
\]
Consider the following commutative diagram:

\phantomsection\label{fac:III:2:no:60:diagram:theorem2}
\[
\begin{tikzcd}[column sep=large,row sep=large]
\mathcal{A}(\Gamma(\mathcal{L}^1))
  \arrow[r] \arrow[d,"\beta"'] &
\mathcal{A}(\Gamma(\mathcal{L}^0))
  \arrow[r] \arrow[d,"\beta"'] &
\mathcal{A}(M) \arrow[r] & 0 \\
\mathcal{L}^1 \arrow[r] &
\mathcal{L}^0 \arrow[r] &
\mathcal{F} \arrow[r] & 0.
\end{tikzcd}
\]
By \hyperref[fac:III:2:no:60:prop:9]{Proposition~9}, the two vertical
homomorphisms are bijective. It follows that $\mathcal{A}(M)$ is
isomorphic to $\mathcal{F}$, as required.
